% Definition file for row def_3_3 -- "EdgeRelations".
% Auto-generated stub by scaffold/solve_chapter.py at row start. The
% manager agent (or a worker dispatched by it) fills the body below with the
% Lean-faithful definition statement(s). Stays close to the lecture notes.
%
% This file is a sibling of `main.tex` inside `<subsection>/tex/`. The
% `subfiles` package lets it compile both standalone (resolving its
% preamble via `main.tex`) and as part of `main.tex` via `\subfile{...}`.

\documentclass[main]{subfiles}

\begin{document}
% ref: def_3_3
% title: EdgeRelations
\def\rowref{def\_3\_3}
\phantomsection\label{def_3_3}

\begin{Def}[EdgeRelations]
    \label{def-edge-relations}
    Let $G = (J, V, E, L)$ be a CDMG (def \ref{def-cdmg}); throughout we use the notation of def \ref{not-cdmg}, and recall that $E \subseteq (J \cup V) \times V$ and $L \subseteq V \times V$, with $L$ symmetric in its two coordinates.

    \begin{enumerate}[label=\roman*.)]
        \item \emph{Adjacency.}
              For $v_1, v_2 \in J \cup V$, the nodes $v_1$ and $v_2$ are called \emph{adjacent in $G$} iff
              \[ (v_1, v_2) \in E \;\lor\; (v_2, v_1) \in E \;\lor\; (v_1, v_2) \in L, \]
              equivalently iff $v_1 \sus v_2 \in G$ (def \ref{not-cdmg}, item~7).

        \item \emph{Edge into a vertex.}
              The relations \emph{into} and \emph{out of} below classify \emph{edges} of $G$ (the underlying ordered pairs of $E$ or of $L$) relative to a single bound vertex parameter $v$ -- not relative to a particular labelling $v_1, v_2$ of the edge's endpoints.
              For an ordered pair $e = (e_1, e_2) \in (J \cup V) \times (J \cup V)$ and a vertex $v \in J \cup V$, we call $e$ an \emph{edge into $v$ (in $G$)} iff
              \[
                \bigl(e \in E \;\land\; e_2 = v\bigr) \;\lor\; \bigl(e \in L \;\land\; (e_1 = v \,\lor\, e_2 = v)\bigr).
              \]
              In particular, only ordered pairs $e \in E \cup L$ can be edges into $v$.

        \item \emph{Edge out of a vertex.}
              For an ordered pair $e = (e_1, e_2) \in (J \cup V) \times (J \cup V)$ and a vertex $v \in J \cup V$, we call $e$ an \emph{edge out of $v$ (in $G$)} iff
              \[ e \in E \;\land\; e_1 = v. \]
              In particular, no $L$-edge is out of any vertex.

        \item \emph{Reconciliation with the source-block writings $v_1 \tuh v_2$, $v_1 \hut v_2$, $v_1 \huh v_2$.}
              The two writings $v_1 \tuh v_2$ and $v_2 \hut v_1$ both abbreviate the single underlying statement $(v_1, v_2) \in E$ (def \ref{not-cdmg}, items~2--3); and the two writings $v_1 \huh v_2$ and $v_2 \huh v_1$ both abbreviate $(v_1, v_2) \in L$, which moreover coincides with $(v_2, v_1) \in L$ by the symmetry of $L$ (def \ref{def-cdmg}).
              Hence membership of an edge in the categories \emph{into $v$} or \emph{out of $v$} depends only on the underlying element of $E$ or of $L$, not on which of these textual writings is used; in particular, the source block's listing $v_1 \tuh v_2$ \emph{or} $v_2 \hut v_1$ in its item~3 and its single-writing listings in item~2 are equivalent under this identification.

              Spelled out per writing, items~ii.--iii.\ above agree with the source block's items~2--3 as follows.
              \begin{itemize}
                  \item An $E$-edge written $v_1 \tuh v_2$ -- i.e.\ underlying ordered pair $(v_1, v_2) \in E$, equivalently written $v_2 \hut v_1$ -- is \emph{into $v_2$} (clause ii.\ fires via $e \in E$ with $e_2 = v_2$) and \emph{out of $v_1$} (clause iii.\ fires via $e \in E$ with $e_1 = v_1$).
                  \item An $E$-edge written $v_1 \hut v_2$ -- i.e.\ underlying ordered pair $(v_2, v_1) \in E$, equivalently written $v_2 \tuh v_1$ -- is \emph{into $v_1$} and \emph{out of $v_2$}, by the same clauses with $v_1$ and $v_2$ interchanged.
                  \item An $L$-edge written $v_1 \huh v_2$ -- i.e.\ underlying ordered pair $(v_1, v_2) \in L$, equivalently $(v_2, v_1) \in L$ and equivalently written $v_2 \huh v_1$ -- is \emph{into $v_1$ and into $v_2$} (the $L$-clause of ii.\ fires for either endpoint), and is \emph{out of neither} $v_1$ nor $v_2$ (clause iii.\ requires $e \in E$, whereas $e \in L$).
              \end{itemize}
              In particular, \emph{into $v$} and \emph{out of $v$} are not complementary classifications of edges incident at $v$: an $L$-edge incident at $v$ is into $v$ but never out of $v$.
    \end{enumerate}
\end{Def}

\end{document}
